% from ega3-7-fr.tex
\oldpage[III]{176}
additif $\rho_*:M\leadsto M_{[\rho]}$
de la catégorie $\boldsymbol{Ab}_B$ des $B$-modules à gauche dans la catégorie
$\boldsymbol{Ab}_A$ des $A$-modules. Par composition, on en déduit un foncteur $T_{(B)}:\boldsymbol{Ab}_B\xrightarrow{\rho_*}\boldsymbol{Ab}_A\xrightarrow{T}\boldsymbol{Ab}$,
évidemment covariant et additif, que nous noterons aussi $T^{(B)}$ (pour des
raisons typographiques) ou $T\otimes_A B$, et que nous dirons obtenu à partir
de $T$ par \emph{extension des scalaires} de $A$ à $B$. Bien entendu, si $B$ est
commutative, on peut considérer $T_{(B)}$ comme un foncteur de
$\boldsymbol{Ab}_B$ dans elle-même (7.1.2). Lorsque $B$ est commutative, et $C$
une $B$-algèbre, on voit aussitôt que $T_{(C)}=(T_{(B)})_{(C)}$;
il est immédiat que l'extension des scalaires est \emph{fonctorielle} et \emph{additive} en
$T$; en outre, lorsque $T$ commute aux limites inductives ou aux sommes
directes (resp. est exact à gauche, exact à droite, exact), il en est de même
de $T_{(B)}$: en effet, $\rho_*$ est exact et commute aux limites inductives
et aux sommes directes.
\end{env}

\begin{env}[7.1.4]
Supposons toujours $A$ commutatif, et soit
$T$ un foncteur covariant additif
$A$-linéaire $\boldsymbol{Ab}_A\to\boldsymbol{Ab}_A$, \emph{commutant aux limites inductives}. Alors, pour toute partie
multiplicative $S$ de $A$ et tout $A$-module $M$, on a un isomorphisme
canonique fonctoriel de $A$-modules
\[
  T(S^{-1}M)\xrightarrow{\sim}S^{-1}T(M).
  \tag{7.1.4.1}\label{III.7.1.4.1.fr}
\]

Supposons en effet d'abord que $S$ soit l'ensemble des puissances $f^n$
$(n\geqslant0)$ d'un élément $f\in A$. On sait alors que $M_f=\varinjlim M_n$,
où $(M_n,\varphi_{nm})$ est le système inductif des $A$-modules $M_n=M$,
avec $\varphi_{nm}:z\leadsto f^{\,n-m}z$ ($\mathbf{0}_{\mathrm I}$, 1.6.1);
d'où dans ce cas l'isomorphisme (7.1.4.1) en vertu de l'hypothèse sur $T$; si
ensuite $S$ est quelconque, $S^{-1}M$ est limite inductive des $M_f$, pour
$f\in S$ ($\mathbf{0}_{\mathrm I}$, 1.4.5) et l'on conclut de la même façon. En outre,
la fonctorialité de l'isomorphisme (7.1.4.1) montre que c'est un isomorphisme
de $S^{-1}A$-\emph{modules}, et qu'on peut donc écrire, à un isomorphisme canonique
près
\[
  T_{(S^{-1}A)}(S^{-1}M)
  =
  S^{-1}T(M)
  =
  T(S^{-1}M).
  \tag{7.1.4.2}\label{III.7.1.4.2.fr}
\]

Lorsque $S=A\setmin\mathfrak p$ est le complémentaire d'un idéal premier
$\mathfrak p$ de $A$, on écrit $T_{\mathfrak p}$ au lieu de
$T_{(A_{\mathfrak p})}$.
\end{env}

\begin{proposition}[7.1.5]
Sous les hypothèses de (7.1.4), si $T_{\mathfrak m}$ est exact à gauche
\textup{(resp.} exact à droite, exact) pour tout idéal maximal $\mathfrak m$ de $A$,
$T$ est exact à gauche \textup{(resp.} exact à droite, exact).
\end{proposition}

\begin{proof}
On sait en effet que lorsque deux sous-modules $N$, $P$ d'un $A$-module $M$
sont tels que $N_{\mathfrak m}=P_{\mathfrak m}$ pour tout idéal maximal
$\mathfrak m$ de $A$, on a $N=P$ (Bourbaki, \emph{Alg. comm.}, chap. II,
\S~3, n\textsuperscript{o}~3, th.~1).
\end{proof}

\subsection{Caractérisations du foncteur produit tensoriel.}
\label{subsection:III.7.2.fr}

\begin{env}[7.2.1]
Soient $A$ un anneau (non nécessairement commutatif), $M$ (resp. $N$) un
$A$-module à gauche (resp. à droite), $P$ un $\mathbf Z$-module. Rappelons
que la donnée d'un $\mathbf Z$-homomorphisme $v:N\otimes_A M\longrightarrow P$
équivaut à la donnée d'une application $\mathbf Z$-bilinéaire $u:N\times M\longrightarrow P$
telle que $u(ta,x)=u(t,ax)$ pour $a\in A$, $t\in N$, $x\in M$,
les deux applications étant reliées par $v(t\otimes x)=u(t,x)$. D'autre part,
la donnée de $u$ équivaut à celle d'un
